\section{News} % (fold)
\label{sec:news}

The documentation you are reading corresponds to the latest version (3.32c) of \tkzNamePack{tkz-elements}.
One significant feature introduced in version 3.0 was the use of the \Iprimitive{directlua} macro, replacing the \tkzNameEnv{tkzelements} environment. If you're creating documents with a large number of figures, it's a good idea to place the \code{init\_elements} function at the start of the lua code, to empty the tables. (see examples in this document)


In  version  3.10c, most functions have been optimized and a few methods are emerging. In particular, methods for determining a circle tangent to different objects. (see \ref{ssub:c_l_pp}; \ref{ssub:method_c__ll__p}; \ref{ssub:method_c__c__pp}; \ref{ssub:method_c_cc_p}; \ref{ssub:method_c_lc_p}; and \ref{ssub:tr_method_c__ll__p})

In version 3.30c, the main new feature is the introduction of the \code{conic} class, which allows you to define parabolas, hyperbolas, and ellipses. The previous \code{ellipse} class is still available but is now obsolete. This new definition is based on a focus, a directrix, and an eccentricity. Curve points are no longer entirely generated using \TIKZ{}; instead, they are computed through a geometric construction, stored in a table, and then passed to \TIKZ{} to draw the curve. This class also includes various methods, such as tangent computation and intersection with a line.

Use the \Imeth{conic}{tkzDrawCoordinates} macro to draw the curve and the \Imeth{conic}{tkzDrawPointOnCurve} macro to place a point on it.

Other important additions in this version include parallel projection onto a line and the affinity transformation.

A small change has also been made to the \code{regular\_polygon} class: the \code{table} attribute has been replaced with the more explicit \code{vertices} attribute.

Another change is the removal of scaling in the \code{lua} part. It turned out that this caused significant code management issues, while scaling with \TIKZ{} was already highly effective (since the most important calculations were already performed). From now on, scaling is not supported by certain functions, so it is recommended to use it only within the \code{tikzpicture} environment.

In version 3.32c, a change has been made to the \code{occs} class: the first argument is still a straight line, but it is no longer the directrix of a conic. Instead, it becomes the main axis passing through a vertex and a focus. This axis will serve as the future ordinate axis. This class can, of course, be used for other constructions beyond conics.

% section news (end)
\endinput

